\section{Arithmetic for the positive determinant}\label{sec:p0-arithmetic}

This section proves the fixed-base denominator estimate for the P0
pencil.  The base is fixed before the limit in $m$ is taken.  No estimate
uniform in unbounded $|a|$ at finite $m$ is asserted.  The proof has an
algebraic part, an exact sum over primes, and a prime-number-theorem limit.
The algebraic and limiting arguments below are all-$m$ arguments.  The
large-ratio contribution is specified by the rational table
\ref{tab:p0-high-potentials} and its type-generation rule.

Write $m=4h$, $g=2h$, $I=\{h,\ldots,5h-1\}$,
$J=\{0,\ldots,8h-1\}$, and $U=8h-1=2m-1$.  Let
\[
 \begin{aligned}
 d_U&=\operatorname{lcm}(1,2,\ldots,U),\\
 N_m&=13h^2-4h=\frac{13}{16}m^2-m,\\
 H_m&=h^2-3h=\frac1{16}m^2-\frac34m.
 \end{aligned}
\]
For a prime $p$, put $d_p=\lfloor\log_p U\rfloor$ and use $v_p$ for the
usual valuation, including its coefficientwise extension to polynomials.
The parameter-dependent entries $A(a),B(a),C(a)$ and determinant
$G_{m,a}(T)$ are those of \eqref{eq:p0det}; in particular the diagonal
conventions in Section~\ref{sec:p0} are retained.

\subsection{The source Mellin identity}
For real $a$ with $|a|>1$ put
\[
 t_i^0(a)=\sum_{r=1}^{i}\frac{a^{i-r}}r,
 \quad b_i^0(a)=\sum_{r=1}^{i}\frac{a^{i-r}}{r^2},
 \quad L(a)=-\sum_{r\ge1}\frac{a^{-r}}r,
 \quad \beta(a)=\sum_{r\ge1}\frac{a^{-r}}{r^2}.
\]
Thus $L(a)=\log(1-1/a)$ and $\beta(a)=\Li_2(1/a)$ with their real
values on both components of $\{|a|>1\}$.  Define
\[
 A^0_{ij}(a)=a^i\delta_{ij},\qquad
 C^0_{ij}(a)=\begin{cases}(a^i-a^j)/(j-i),&i\ne j,\\0,&i=j,\end{cases}
\]
and
\[
 B^0_{ij}(a)=\begin{cases}(t_i^0(a)-t_j^0(a))/(j-i),&i\ne j,\\
 b_i^0(a),&i=j.\end{cases}
\]

\begin{lemma}[source Mellin identity]
\label{lem:p0-source-mellin}
For $|a|>1$ and all $i,j\ge0$,
\begin{equation}\label{eq:p0-source-mellin}
 M_{ij}(a):=B^0_{ij}(a)+L(a)C^0_{ij}(a)-\beta(a)A^0_{ij}(a)
 =-\sum_{r\ge1}\frac{a^{-r}}{(i+r)(j+r)}
 =-\int_0^1\!\!\int_0^1\frac{x^iy^j}{a-xy}\,dx\,dy.
\end{equation}
\end{lemma}
\begin{proof}
The two finite identities
\[
 t_i^0(a)+a^iL(a)=-\sum_{r\ge1}\frac{a^{-r}}{i+r},\qquad
 b_i^0(a)-a^i\beta(a)=-\sum_{r\ge1}\frac{a^{-r}}{(i+r)^2}
\]
follow by separating the first $i$ terms of the defining series.  Taking the
divided difference of the first identity and using
\[
 \frac{(i+r)^{-1}-(j+r)^{-1}}{j-i}
 =\frac1{(i+r)(j+r)}
\]
proves the off-diagonal case of \eqref{eq:p0-source-mellin}; the second
identity proves the diagonal case.  Since $|a|>1$, the geometric series
$(a-xy)^{-1}=\sum_{r\ge1}a^{-r}(xy)^{r-1}$ converges uniformly on
$[0,1]^2$, so termwise integration gives the last equality.
\end{proof}

Put
\[
 D_g f(a)=\frac1{g!}\frac{d^g}{da^g}\bigl(a^gf(a)\bigr).
\]
On every compact subset of $\{|a|>1\}$ the series in
\eqref{eq:p0-source-mellin}, after multiplication by $a^g$, and all its
derivatives through order $g$, are dominated by $C(1+r)^g\rho^r$ with
$\rho<1$.  Termwise differentiation is therefore valid and
\begin{equation}\label{eq:p0-source-Dg}
 D_gM_{ij}(a)=(-1)^{g+1}\sum_{r>g}
 \frac{\binom{r-1}{g}a^{-r}}{(i+r)(j+r)}.
\end{equation}
Indeed $D_ga^{-r}=0$ for $r\le g$ and equals
$(-1)^g\binom{r-1}{g}a^{-r}$ for $r>g$.

For completeness, the right side is exactly the finite rational table used in
the determinant.  If $U_i(a)=a^i\binom{i+g}{g}$ and
$Q_s(a)=(a/(a-1))^s-1$, the finite definitions of $A,B,C$ in
Section~\ref{sec:p0} already use $a$ throughout.  Leibniz's rule gives
$D_ga^i=U_i(a)$ and, for $s\ge1$,
\[
 \frac{a^sL^{(s)}(a)}{s!}=\frac{(-1)^{s-1}}sQ_s(a),\qquad
 \frac{a^s\beta^{(s)}(a)}{s!}=\frac{(-1)^{s-1}}s
 \left(L(a)-\sum_{r=1}^{s-1}\frac{Q_r(a)}r\right).
\]
The latter follows from $\beta'(a)=L(a)/a$.  Substitution in
\eqref{eq:p0-source-mellin} proves
\begin{equation}\label{eq:p0-source-identity}
 D_gM_{ij}(a)=B^{(g)}_{ij}(a)+L(a)C^{(g)}_{ij}(a)
                         -\beta(a)A^{(g)}_{ij}(a).
\end{equation}
For $i=j$, the coefficient of $L(a)$ is
$-U_i(a)\sum_{s=1}^g(i+s)^{-1}$, because
\[
 \sum_{s=1}^g\frac{(-1)^{s-1}}s\binom{i+g}{g-s}
 =\binom{i+g}{g}\sum_{s=1}^g\frac1{i+s}.
\]
This identity is the derivative at $x=0$ of
$\binom{i+g+x}{g}=\sum_{s=0}^g\binom{i+g}{g-s}\binom{x}{s}$;
the coefficient of $-\beta$ is $U_i(a)$.  Hence the diagonal conventions,
and not only the off-diagonal divided differences, are covered by
\eqref{eq:p0-source-identity}.

\subsection{The polynomial divisor}

We first record the two polynomial divisibilities that account for the
exceptional primes.  All statements in this subsection are identities for
every $h\ge1$.

\begin{lemma}[formal divisor at $a=0$]\label{lem:p0-a-divisor}
With $b=a-1$,
\begin{equation}\label{eq:p0-denom}
 d_U^{3m}b^{gm}G_{m,a}(T)\in\mathbb Z[a,T],
 \qquad
 d_U^{3m}b^{gm}G_{m,a}(T)\in a^{N_m}\mathbb Z[a,T].
\end{equation}
\end{lemma}

\begin{proof}
The only parameter denominators in the finite source formulas are powers of
$b$ in
\[
 Q_s(a)=\left(\frac{a}{a-1}\right)^s-1.
\]
After multiplying each lower row by $b^g$, its occurrence becomes the integer
polynomial
\[
 R_s(a)=b^gQ_s(a)=a^s b^{g-s}-b^g\in\mathbb Z[a]
 \quad(1\le s\le g).
\]
Its constant coefficient is $-(-1)^g$, a unit.  The other source
denominators are harmonic indices, summation indices, and divided-difference
denominators.  Every such positive index is at most $U$: for example
$i+g\le7h-1<U$, and $|j-i|\le7h-1<U$.  Multiplication of each top row by
$d_U$ and each lower row by $d_U^2$ clears them.  Thus
\[
 P_m(a,T):=d_U^{3m}b^{gm}G_{m,a}(T)\in\mathbb Z[a,T].
\]

To find its order at $a=0$, work in $\mathbb Q[[a]]$.  Set
\[
 c_i=\binom{i+g}{g},\quad u_i=c_i a^i,\quad
 h_i=\sum_{s=1}^g\frac1{i+s},\quad
 \gamma_i=\frac12\left(h_i^2-\sum_{s=1}^g\frac1{(i+s)^2}\right),
\]
and $L_1=\sum_{r\ge1}a^r/r$, $L_2=\sum_{r\ge1}a^r/r^2$.  The source
identities, coefficientwise in $a$, are
\begin{align}
 t_i-L_1u_i&=\sum_{k\ne i}\frac{c_ka^k}{i-k}-u_i h_i,
 \label{eq:p0-formal-t}\\
 b_i+L_1u_i h_i+L_2u_i
 &=\sum_{k\ne i}\frac{c_ka^k}{(i-k)^2}+\gamma_i u_i.
 \label{eq:p0-formal-b}
\end{align}
Here the $b_i$ on the left is the diagonal source entry, not the parameter
$b=a-1$.  We use a direct coefficient check, so no auxiliary logarithmic
symbol is being differentiated.  For $r\ge0$, expand every finite source
entry in $\mathbb Q[[a]]$, replace
\[
 [a^r]L_1=\frac1r\quad(r\ge1),\qquad
 [a^r]L_2=\frac1{r^2},
\]
and use
\[
 [a^r]u_i=\binom{i+g}{g}[r=i],\qquad
 \gamma_i=\sum_{1\le s<t\le g}\frac1{(i+s)(i+t)}.
\]
The coefficient of the first source identity is then
\[
 [a^r]t_i-[a^r](L_1u_i)=
 \begin{cases}c_r/(i-r),&r\ne i,\\-c_i h_i,&r=i,\end{cases}
\]
which is exactly the right side of \eqref{eq:p0-formal-t}.  For the second
identity, the coefficient is
\[
 [a^r]\bigl(b_i+L_1u_i h_i+L_2u_i\bigr)=
 \begin{cases}c_r/(i-r)^2,&r\ne i,\\c_i\gamma_i,&r=i.\end{cases}
\]
This is the coefficient of its right side, including the diagonal correction.
All sums are finite at a fixed coefficient, so these are identities in
$\mathbb Q[[a]]$ coefficient by coefficient.  The only binomial identity
used in collecting the harmonic terms is the differentiated finite identity
\[
 \binom{i+g+x}{g}=\sum_{s=0}^g\binom{i+g}{g-s}\binom{x}{s}
\]
at $x=0$; hence the argument does not assume a formal primitive for $1/a$.

For an independent indeterminate $X$, form
\[
 \mathcal H(a,X)=
 \begin{bmatrix}
 C(a)\\ B(a)-L_1(a)C(a)+L_2(a)A(a)-XA(a)
 \end{bmatrix}.
\]
For a fixed $k\ge0$, let $e_I(i)=[i=k]$,
$v_I(i)=(i-k)^{-1}$ for $i\ne k$ and $v_I(k)=0$, and
$w_I=v_I^2$; define $e_J,v_J,w_J$ in the same way.  Equations
\eqref{eq:p0-formal-t}--\eqref{eq:p0-formal-b} give the exact rank-one
expansion
\begin{equation}\label{eq:p0-rank-one}
 [a^k]\mathcal H=c_k\left(
 \binom{e_I}{v_I-h_ke_I}v_J^{\mathsf T}
 +\binom{0}{-e_I}w_J^{\mathsf T}
 +\binom{v_I-h_ke_I}{-w_I-h_kv_I+(\gamma_k-X)e_I}e_J^{\mathsf T}
 \right).
\end{equation}
In entries, the lower block of this coefficient is
\[
 [a^k]\bigl(B-L_1C+L_2A-XA\bigr)_{ij}
 =c_k\begin{cases}
 ((i-k)(j-k))^{-1},&i\ne k,\ j\ne k,\\
 -h_k(j-k)^{-1}-(j-k)^{-2},&i=k,\ j\ne k,\\
 -(i-k)^{-2}-h_k(i-k)^{-1},&i\ne k,\ j=k,\\
 \gamma_k-X,&i=j=k.
 \end{cases}
\]
The first feature has weight $k$; the second is present only when $k\in I$,
and the third only when $k\in J$.  A nonzero $2m$-row determinant term must
choose $2m$ distinct rank-one features.  Since $I=[h,5h)$ and $J=[0,8h)$,
the least possible total feature weight is obtained from two copies of
$0,\ldots,h-1$ and three copies of $h,\ldots,3h-1$.  Its value is
\[
 2\sum_{k=0}^{h-1}k+3\sum_{k=h}^{3h-1}k=13h^2-4h=N_m.
\]
Finite multilinearity is legitimate coefficientwise: a coefficient of $a^r$
can involve only features of weight at most $r$.  Thus
$a^{N_m}\mid\det\mathcal H(a,X)$ in $\mathbb Q[[a]][X]$.  The determinant-one
row shear gives
$\det\mathcal H(a,X)=G_{m,a}(X-L_2(a))$; since $L_2\in a\mathbb Q[[a]]$,
translation preserves the order at zero.  Therefore $a^{N_m}$ divides
$P_m$ in $\mathbb Q[a,T]$.  As $P_m\in\mathbb Z[a,T]$ and the divisor is
the monomial $a^{N_m}$, the quotient is also in $\mathbb Z[a,T]$.
\end{proof}

\begin{lemma}[endpoint pole divisor]\label{lem:p0-pole}
Coefficientwise in $T$,
\[
 \operatorname{ord}_{a=1}G_{4h,a}(T)\ge3h-h^2.
\]
Consequently
\begin{equation}\label{eq:p0-refined-divisor}
 d_U^{3m}(a-1)^{H_m}G_{m,a}(T)
       \in a^{N_m}\mathbb Z[a,T].
\end{equation}
The exponent $3h-h^2$ is a sufficient lower bound; equality for every $h$ is
not asserted.
\end{lemma}

\begin{proof}
Put $b=a-1$ and $f(x)=\binom{x+g}{g}$.  At $a=1$, the full top block,
including its diagonal, is the divided-difference polynomial
\[
 C(x,y)=\frac{f(x)-f(y)}{y-x},\qquad C(x,x)=-f'(x),
\]
of total degree at most $g-1$.  For the lower block, the source Mellin
identity gives
\[
 (D_gM_{ij})(a)=(-1)^{g+1}\sum_{r\ge1}
   \binom{r-1}{g}\frac{a^{-r}}{(r+i)(r+j)}
 =B_{ij}+L(a)C_{ij}-\Li_2(1/a)A_{ij}.
\]
Polynomial division in $r$ yields
\[
 \frac{\binom{r-1}{g}}{(r+i)(r+j)}
 =\sum_{d=0}^{g-2}q_d(i,j)r^d+O(r^{-1}),
 \qquad \deg_{i,j}q_d\le g-2-d.
\]
The power sum $\sum_{r\ge1}r^d(1+b)^{-r}$ has pole order at most $d+1$;
the remainder is $O(1+|\log b|)$.  Since $B$ is rational in $b$, while
$LC-\Li_2(1/a)A=O(1+|\log b|)$, this identifies its entire principal
part, including diagonal entries:
\[
 [b^{-s}]B_{ij}(1+b)=P_s(i,j),\qquad
 \deg P_s\le g-1-s\quad(1\le s\le g-1).
\]
Apply forward differences of orders $u=0,\ldots,m-1$ to the rows in each
block and $v=0,\ldots,2m-1$ to columns.  These are unitriangular transforms,
so they preserve the determinant.  The degree bounds imply
\[
 \operatorname{ord}_b\widehat C_{uv}\ge[ u+v\ge g],\qquad
 \operatorname{ord}_b(\widehat B-T\widehat A)_{uv}
       \ge\min(0,u+v-g+1).
\]
In a determinant term, let $k$ be the number of lower entries with a negative
bound.  Then $0\le k\le g-1$, and their distinct row and column indices
give a total order at least $k(k-g)$.  Their columns lie among
$0,\ldots,g-2$; hence at most $g-k$ top entries can have order-zero bounds,
so at least $m-g+k$ top entries contribute order at least one.  Every
determinant term has order at least
\[
 m-g+k+k(k-g)=m-g+k(k-g+1).
\]
For $m=4h$, $g=2h$, this is
\[
 2h+(k-h)(k-h+1)-h(h-1)\ge3h-h^2,
\]
because $(k-h)(k-h+1)\ge0$ for integral $k$.  This proves the order bound,
including all diagonal entries.

By Lemma~\ref{lem:p0-a-divisor},
$P_m=d_U^{3m}b^{8h^2}G=a^{N_m}S$ with $S\in\mathbb Z[a,T]$.  The pole
bound shows that $b^{7h^2+3h}$ divides $P_m$ in $\mathbb Q[a,T]$.
This divisor is monic, so division preserves integrality; it is coprime to
$a^{N_m}$.  Cancelling it yields
\[
 d_U^{3m}b^{8h^2-(7h^2+3h)}G
 =d_U^{3m}b^{H_m}G\in a^{N_m}\mathbb Z[a,T],
\]
as claimed.  The argument also covers $h=1,2$, when $H_m=-2$; in those
cases the displayed expression is an identity of rational functions whose
right side is the stated integer polynomial.
\end{proof}

\subsection{Primewise valuation bounds}

For a fixed integer $a$ with $|a|\ge2$, define
\[
 F_{m,p}(a)=3m d_p+H_m v_p(a-1)-N_m v_p(a).
\]
We next define a second, moving-prime bound.  For an odd prime satisfying
$p^2>U+g$, put
\[
 r=g\bmod p,\qquad
 \kappa_n=[n\bmod p\ge p-r],\qquad
 \theta_n=(\kappa_n-1)[n\ge p\text{ or }p\le g],
\]
where $\kappa_n=0$ if $r=0$.  Define the integer $W_m(p)$ by the following
deterministic rule (it is the chosen upper bound, not an undefined optimum):
\begin{itemize}
\item if $p^2\le U+g$, set $W_m(p)=3m d_p$;
\item if $p^2>U+g$ and $p$ lies on a positive-ratio wall, set $W_m(p)=3m$;
\item otherwise, for $p/m<1/2$, set $W_m(p)$ to the smaller of the two
  integer right sides in \eqref{eq:p0-low-W}; for
  $1/2<p/m<5/2$, generate the groups from the affine endpoints below,
  sum the displayed row and column potentials in Table~\ref{tab:p0-high-potentials}
  over those groups, and take that exact integer objective;
\item if $p>U+g$, set $W_m(p)=0$.
\end{itemize}
At the low-ratio endpoint $p/m=1/2$, the wall convention $3m$ applies.
The group widths are integers because $4\mid m$ and all endpoint intercepts
have denominator at most four; consequently the high-ratio objective is an
integer, not a real relaxation.  These clauses define $W_m(p)$ for every
prime before it is used in \ref{lem:p0-entry-bounds}.
The finite collection of walls in $[1/128,5/2]$ is defined below.  These
rules are compatible at overlapping cases: the generic denominator bound is
used first, and $p>U+g$ is outside every positive-ratio cell.

\begin{lemma}[coefficientwise entry bounds]\label{lem:p0-entry-bounds}
For every integer specialization $a$ and every coefficient in $T$, the
valuation of the determinant satisfies
\[
 -v_p(G_{m,a})\le W_m(p)+gm\,v_p(a-1).
\]
Also, for every prime $p$,
\[
 -v_p(G_{m,a})\le F_{m,p}(a).
\]
\end{lemma}

\begin{proof}
The second inequality follows directly from
\eqref{eq:p0-refined-divisor}.  For the first, multiply every lower row by
$b^g$, so that
$\det[C;b^g(B-TA)]=b^{gm}G_{m,a}(T)$.  When $p^2\le U+g$, the factors
$d_U$ and $d_U^2$ clear top and lower rows, giving the claimed $3md_p$
bound.  When $p>U+g$, every source denominator is a $p$-unit and the scaled
entries are $p$-integral.

It remains to consider odd $p$ with $p^2>U+g$.  Every relevant upper index
is less than $p^2$.  Kummer's carry rule therefore gives
\[
 v_p\binom{n+g}{g}=\kappa_n.
\]
If $p\mid s$, then $g-s$ has the same units digit as $g$, so
\[
 v_p\binom{n+g}{g-s}\ge\kappa_n;
\]
 and, if $p\mid s$ and $s\le n$,
$v_p\binom{n-s+g}{g}\ge\kappa_n$.  No second carry is possible because
all sums are below $p^2$.  After multiplying by $b^g$, each $Q_s$ is the
primitive polynomial $R_s$ above, whose Gauss valuation is zero.  In these
sums, terms with $p\nmid s$ are $p$-integral; in the mixed
$B_{ii}$ terms, a single $p$-denominator contributes at worst
$-1\ge\kappa_i-2$, since $\kappa_i\le1$.  Every nonzero source denominator
and every row or column difference is less than
$p^2$, hence has $p$-valuation at most one.  The ultrametric inequality now
gives the following lower bounds for scaled entries:
\begin{center}
\renewcommand{\arraystretch}{1.12}
\normalsize
\begin{tabular}{@{}p{0.43\linewidth}p{0.46\linewidth}@{}}
\toprule
Entry & Lower bound for $v_p$ \\
\midrule
$C_{ij}$, $i\ne j$ & $\min(\kappa_i,\kappa_j)-[i\equiv j\pmod p]$ \\
$C_{ii}$ & $(\kappa_i-1)[p\le g]$ \\
$(b^gB)_{ij}$, $i\ne j$ & $\min(\theta_i,\theta_j)-[i\equiv j\pmod p]$ \\
$(b^gB)_{ii}$, $i\ge p$ or $2p\le g$ & $\kappa_i-2$ \\
$(b^gB)_{ii}$, $i<p$ and $p\le g<2p$ & $-1$ \\
$(b^gB)_{ii}$, $i<p$ and $p>g$ & $0$ \\
\bottomrule
\end{tabular}
\end{center}
For $C_{ii}$, the case $p\le g$ uses the general bound
$v_p(U_i/(i+s))\ge\kappa_i-1$, since $i+s<p^2$.  If $p>g$ and
$p\mid i+s$, then $1\le s\le g<p$, so $i\bmod p=p-s$; adding $g$ to
$i$ produces exactly one base-$p$ carry and hence $\kappa_i=1$.  Thus every
term with a $p$-divisible denominator is still $p$-integral, and all other
terms are $p$-integral because $v_p(U_i)=\kappa_i\ge0$.  For the lower diagonal, a squared denominator
in the first sum requires $i\ge p$; two $p$-denominators in the correction
sum require $2p\le g$.  In either case the relevant binomial coefficient
has valuation at least $\kappa_i$.  If $p\le g<2p$ and $i<p$, at most one
denominator is divisible by $p$; if $p>g$ and $i<p$, none is.  The term
$-TA$ is integral after the row scaling and does not weaken these bounds.

By the Leibniz expansion, minus the determinant valuation is at most the
maximum assignment weight formed by negating these entry lower bounds.
The exact profile $W_m(p)$ bounds that assignment in each range.  The
certificate details follow.
\end{proof}

\subsection{Exact moving-prime profile}

Let $\rho=p/m$.  For a residue $v\in\{0,\ldots,p-1\}$, the copy counts are
\[
 c_I(v)=\#\{n\in\mathbb Z:1/4\le v/m+n\rho<5/4\},\qquad
 c_J(v)=\#\{n\in\mathbb Z:0\le v/m+n\rho<2\}.
\]
The carry flag is constant on each residue group.  For $p<g$, let $K_I,K_J$
count the killed row and column copies ($\kappa_v=1$), and put
\[
 S_U=\sum_{\kappa_v=0}\min(2c_I(v),c_J(v)),\qquad
 S_K=\sum_{\kappa_v=1}\min(2c_I(v),c_J(v)).
\]
The two feasible assignment bounds are
\begin{equation}\label{eq:p0-low-W}
 W_m(p)\le\min\{m+S_U,\;3m-2K_I-K_J+S_K\}.
\end{equation}
Indeed, after subtracting the lower-row baseline $m$, the assignment score is
$x-z$, where $x$ counts unkilled same-residue matches and $z$ killed
off-residue matches.  If $y$ is the number of killed same-residue matches,
then $x\le S_U$, $y\le S_K$, and counting killed row and column copies gives
$z+y\ge2K_I+K_J-2m+x$.  These yield the two inequalities.

All group endpoints are generated by a finite exact rule.  For a fixed open
wall cell and any interior rational sample $\rho_0$, set
$q=\lfloor1/(2\rho_0)\rfloor$ and order the affine endpoints
\[
 (0,0),\ (0,1),\ (-1/2,q+1),\quad
 (x,-\lfloor x/\rho_0\rfloor)\ (x=1/4,5/4,2),
\]
by their values $A+B\rho_0$.  Consecutive endpoints give a group of width
$A+B\rho$.  At its midpoint $t$, the $I$- and $J$-copy lists are respectively
\[
 \left\lceil\frac{1/4-t}{\rho_0}\right\rceil\le n<
 \left\lceil\frac{5/4-t}{\rho_0}\right\rceil,
 \qquad
 \left\lceil\frac{-t}{\rho_0}\right\rceil\le n<
 \left\lceil\frac{2-t}{\rho_0}\right\rceil,
\]
and the group is killed precisely when
$t\ge(q+1)\rho_0-1/2$.  These data remain constant throughout the open cell.
The full wall set in a compact ratio interval is generated by
\[
\mathcal W([r_0,r_1])=\{r_0,r_1\}\cup
\left\{\frac{x-y}{n}:\begin{aligned}[t]
&x,y\in\{0,1/4,5/4,2,-1/2\},\quad x>y,\\
&n\in\mathbb N,\quad r_0<\frac{x-y}{n}<r_1
\end{aligned}\right\}.
\]
It is finite because $x-y\le5/2$ and $n<(x-y)/r_0$.

For $\rho\in[\delta,1/2]$, $\delta=1/128$, the rule gives 756 lower-envelope
affine pieces.  Their exact integral is
\begin{equation}\label{eq:p0-low-integral}
 L_*=
 \frac{\begin{array}{c}
 \texttt{1448276314740194964855260100800462127260298}\\
 \texttt{03104279674952972562737001039673489563365828533335}\\
 \texttt{0450905704443029571002838710237}
 \end{array}}
 {\begin{array}{c}
 \texttt{140409888158080315725937691852162573583946167555}\\
 \texttt{38734362972176526622161627829570435939660600468548}\\
 \texttt{26150978176673015236224000}
 \end{array}}.
\end{equation}
In each array, concatenate its three rows to obtain the integer.  This value
is obtained by integrating the minimum in \eqref{eq:p0-low-W}
on each rational group cell, splitting at its rational affine crossings.
The formula above is the exact reconstruction rule; no optimizer is used.

For $0<\rho<\delta$, zero column potentials give
$W_m(p)\le3m-2K_I$.  Since an interval of $m$ consecutive integers contains
at least $\lfloor m/p\rfloor(g\bmod p)$ killed indices, the profile
$f_{\rm tail}(\rho)=3-2\{1/(2\rho)\}+2\rho$ satisfies
$W_m(p)/m\le f_{\rm tail}(\rho)$.
Writing $\delta=1/(2Q)$ with $Q=64$ and substituting
$x=1/(2\rho)$ gives
\[
 \int_0^\delta\left(3-2\left\{\frac1{2\rho}\right\}\right)d\rho
 =\frac1{2Q}+\sum_{q\ge Q}\left(\frac1q-\log(1+1/q)\right).
\]
For each $q\ge Q$, the summand differs from $1/(2q(q+1))$ by
\[
 \int_0^1\frac{s(1-s)}{q(q+s)(q+1)}\,ds
 \le\frac1{6q^2(q+1)}.
\]
Summing this bound and integrating the added $2\rho$ term yields
\begin{equation}\label{eq:p0-small-tail}
 \int_0^\delta f_{\rm tail}(\rho)\,d\rho\le
 2\delta+2\delta^2=\frac{129}{8192}.
\end{equation}

For $1/2<\rho<5/2$, $p>g$, so each group has a carry bit $k\in\{0,1\}$.
If $n$ denotes a copy index, put
$\theta=(k-1)[n\ge1]$.  A row type is $(j,n,e,k)$, where $j$ is its group,
$n$ is an $I$-copy and $e\in\{0,1\}$ is the top or lower block; a column
type is $(j,n,k)$ with $n$ a $J$-copy.  Order row types by increasing group,
then copy, then $e=0,1$; order column types by increasing group then copy.
For a row type $r$ and column type $c$, the required potential inequality is
\[
 u_r+v_c\ge w_{rc}.
\]
For different groups the right side is $-\min(k_r,k_c)$ in the top block
and $-\min(\theta_r,\theta_c)$ in the lower block.  For the same group but
distinct copies add $1$.  For a same-index top diagonal use $0$; for a
same-index lower diagonal use $(2-k_r)$ if $n\ge1$, and $0$ if $n=0$.
These are exactly the negated entry bounds of
Lemma~\ref{lem:p0-entry-bounds}, including both diagonal cases.  The
following integer vectors give feasible potentials on every large-ratio
interval; spaces separate successive coordinates and the semicolon separates
the row and column vectors.  The group and copy ordering is the rule just
given.  Each group contains exactly $m(A_j+B_j\rho)$ residues; this is an
integer because $m$ is divisible by four and all endpoint intercepts have
denominator at most four.  Thus the potential objective is an integer upper
bound at every prime ratio, not merely a real-valued relaxation.

\renewcommand{\arraystretch}{1.15}
\normalsize
\begin{longtable}{@{}c p{0.33\linewidth}p{0.34\linewidth}c@{}}
\caption{Exact large-ratio assignment potentials.  The last column is
  $\sum_r u_r(A_{j(r)}+B_{j(r)}\rho)+
    \sum_c v_c(A_{j(c)}+B_{j(c)}\rho)$.}\label{tab:p0-high-potentials}\\
\toprule
Ratio cell & Row potentials $u$ & Column potentials $v$ & Objective \\
\midrule
\endfirsthead
\toprule
Ratio cell & Row potentials $u$ & Column potentials $v$ & Objective \\
\midrule
\endhead
$(1/2,7/12)$ &
\texttt{1 2 1 2 0 1 0 1 1 2 0 1 0 1 0 1 0 1} &
\texttt{0 0 0 0 0 0 0 0 -1 -1 -1 -1 0 0 0 0 0 0 0} & $1$ \\
$(7/12,5/8)$ &
\texttt{1 2 1 2 1 2 1 2 1 2 0 1 0 1} &
\texttt{0 0 0 0 0 0 0 0 -1 -1 -1 -1 -1 -1 -1 0 0 0} & $3\rho-3/4$ \\
$(5/8,2/3)$ &
\texttt{1 2 1 2 1 2 0 1 0 1 1 2} &
\texttt{0 0 0 0 0 0 0 -1 -1 -1 0 0 0 -1 -1 -1} & $\rho+1/2$ \\
$(2/3,3/4)$ &
\texttt{1 2 1 2 0 1 0 1 1 2 0 1} &
\texttt{0 0 0 -1 -1 -1 0 0 0 -1 -1 -1 0 0} & $4\rho-3/2$ \\
$(3/4,5/6)$ &
\texttt{1 2 0 1 0 1 0 1 0 1 0 1 0 1 0 1} &
\texttt{0 0 0 1 1 1 0 0 0 0 0 0 0} & $3\rho-3/4$ \\
$(5/6,7/8)$ &
\texttt{1 2 0 1 0 1 0 1 0 1 0 1 0 1 0 1} &
\texttt{0 0 0 1 1 1 1 1 0 0 0 0} & $7/4$ \\
$(7/8,1)$ &
\texttt{1 2 0 1 0 1 0 1 0 1 0 1} &
\texttt{0 0 0 1 1 1 1 0 1 0 0} & $7/4$ \\
$(1,5/4)$ &
\texttt{1 2 1 1 1 1 1 1} &
\texttt{0 0 -1 0 -1 1 -1 0 -1} & $15/4-2\rho$ \\
$(5/4,7/4)$ &
\texttt{0 0 0 0 0 0} &
\texttt{0 1 0 2 0 0 0} & $15/4-2\rho$ \\
$(7/4,2)$ &
\texttt{0 0} &
\texttt{0 1 0 0 0 0} & $2-\rho$ \\
$(2,5/2)$ & $\texttt{0 0}$ & $\texttt{0 0 0 0}$ & $0$ \\
\bottomrule
\end{longtable}

The eleven large-ratio intervals are exactly the walls generated by
$\mathcal W([1/2,5/2])$.  For each interval, the group formula above
determines the row and column types, and the displayed vectors satisfy all
the pairwise rational inequalities.  The edge, diagonal, and wall comparisons
are generated by the same displayed rule.  Summing the
type potentials times their affine group widths
gives the objective column; taking its exact integral over the high cells
yields
\[
 \int_{1/2}^{5/2}f_{\rm high}(\rho)\,d\rho=\frac{287}{192}.
\]
At a wall use $W_m(p)=3m$; there are finitely many walls, so this changes
the prime sum by $O(m\log m)$ for each $m$.

The exact profile constant is
\begin{align}
 K&=L_*+\frac{129}{8192}+\frac{287}{192}\\
 &=\frac{\begin{array}{c}
 \texttt{28553776527336974955476146348812001767941}\\
 \texttt{86855224421109769015232357578975297803041084792395}\\
 \texttt{8183617568292693602659280722163521}
 \end{array}}
 {\begin{array}{c}
 \texttt{112327910526464252580750153481730058867156934044}\\
 \texttt{30987490377741221297729302263656348751728480374838}\\
 \texttt{609207825413384121889792000}
 \end{array}}\\
 &=2.542001929307655\ldots<2.542002.
\label{eq:p0-K}
\end{align}
The displayed rational table and the type-generation rule are the complete
finite input for this part of the estimate.

\subsection{Exceptional primes and the PNT limit}

Define the rational scalar
\begin{equation}\label{eq:p0-delta}
 \Delta_{m,a}=\prod_p p^{E_{m,p}(a)},\qquad
 E_{m,p}(a)=\min\{W_m(p)+gm\,v_p(a-1),\ F_{m,p}(a)\}.
\end{equation}
Negative exponents are allowed; $\Delta_{m,a}>0$ is a rational number.
The product is finite: if $p>U+g$ and $p\nmid a(a-1)$, both entries in
the minimum are zero.  Lemma~\ref{lem:p0-entry-bounds} implies
coefficientwise that
\begin{equation}\label{eq:p0-integerization}
 \Delta_{m,a}G_{m,a}(T)\in\mathbb Z[T].
\end{equation}
This is a specialization-wise p-adic argument; it does not claim that the
sharp scalar $\Delta_{m,a}$ is the specialization of an integer polynomial
in $a$.

Since $\gcd(a,a-1)=1$, the exceptional sets $p\mid a$ and $p\mid a-1$
are disjoint.  For a fixed $p\mid a$,
\[
 F_{m,p}=O(m\log m)-N_m v_p(a),
 \qquad W_m(p)+gmv_p(a-1)=W_m(p),
\]
and the quadratic negative term makes $F_{m,p}$ the minimum for all
sufficiently large $m$.  For a fixed $p\mid a-1$,
\[
 F_{m,p}=O(m\log m)+H_m v_p(a-1),\qquad
 W_m(p)+gmv_p(a-1),
\]
and $gm-H_m=7h^2+3h$ makes $F_{m,p}$ eventually the minimum.  These
comparisons apply also to $p=2$, which divides exactly one of $a,a-1$.
For every other fixed prime the moving profile applies, with the
$p^2\le U+g$ range contributing $o(m^2)$ in aggregate:
Put $X=U+g$.  The squarefull part is bounded without counting a
prime more than once by
\[
 \sum_{p\le\sqrt X}3m d_p\log p
 \le 3m\sum_{j=1}^{\lfloor\log_2 U\rfloor}
    \vartheta\!\left(\min\{\sqrt X,U^{1/j}\}\right)
 =O(m^{3/2}\log m)=o(m^2),
\]
where $\vartheta(x)=\sum_{p\le x}\log p$.  The displayed equality is an
upper-bound comparison (the left side is bounded by the right side),
since $d_p$ is the number of admissible exponent layers and each layer is
charged to $p^j\le U$; the extra $\log m$ is harmless after normalization.

On any fixed positive ratio interval, $W_m(p)\le m f(p/m)$ with $f$ bounded
and piecewise affine, apart from the finite wall values.  The weighted prime
number theorem, obtained from the ordinary PNT by partial summation, gives
\[
 \frac1m\sum_p f(p/m)\log p\longrightarrow\int f(\rho)\,d\rho.
\]
For $0<\rho<\eps$, first remove the already negligible range
$p\le\sqrt{U+g}$ and then omit $\sqrt{U+g}<p\le\eps m$; there $d_p=1$ and
the crude bound $3m\vartheta(\eps m)$ has normalized limsup at most
$3\eps$.  Apply the weighted PNT on $[\eps,\delta]$ and then let
$\eps\downarrow0$.  The exact
tail bound \eqref{eq:p0-small-tail}, the low integral \eqref{eq:p0-low-integral},
and the high integral therefore imply
\[
 \limsup_{\substack{m\to\infty\\4\mid m}}
 \frac1{m^2}\sum_{p\nmid a(a-1)}E_{m,p}(a)\log p\le K.
\]
The fixed exceptional primes contribute
\[
 \frac{H_m}{m^2}\log|a-1|
 -\frac{N_m}{m^2}\log|a|
 \longrightarrow
 \frac1{16}\log|a-1|-\frac{13}{16}\log|a|.
\]
We have proved the following fixed-base estimate.

\begin{proposition}[P0 fixed-base denominator rate]\label{prop:p0-sharp-arith}
For every fixed integer $|a|\ge2$, the rational number $\Delta_{m,a}$ in
\eqref{eq:p0-delta} satisfies \eqref{eq:p0-integerization} and
\[
 \limsup_{\substack{m\to\infty\\4\mid m}}
 \frac{\log\Delta_{m,a}}{m^2}
 \le K+\frac1{16}\log|a-1|-
             \frac{13}{16}\log|a|,
 \qquad K<2.542002.
\]
The statement is for each fixed base.  In particular, its finite-$m$
remainder is not uniform as $|a|\to\infty$.
\end{proposition}

\paragraph{Completion of the estimate.}
The formal divisor and pole lemmas, Kummer bounds, residue-cell formulas,
exceptional-prime comparisons, and PNT passage are all-$m$ derivations.
The low profile is reconstructed by the rational wall-and-crossing rule
above.  The high profile is hand-checkable from the printed potential vectors
and the stated type-generation rule: each vector pair reduces the full
assignment inequality to finitely many integer comparisons.  The low profile
has 756 affine pieces and the high profile eleven cells.  The all-$m$
conclusion uses the displayed identities, the finite comparison count above,
and the weighted prime number theorem.
